\documentclass[Afour,sagev,times]{sagej}

% Optional packages commonly used with SAGE class
\usepackage{url}
\usepackage[colorlinks=true,
            linkcolor=blue,
            urlcolor=blue,
            citecolor=blue]{hyperref}
\hypersetup{colorlinks=true,linkcolor=black,citecolor=black,urlcolor=black}
\usepackage{graphicx}
\usepackage{amsmath}
\usepackage{enumitem}
%\usepackage[numbers,super]{natbib}

\begin{document}

\runninghead{Short Running Title}
\title{Full Manuscript Title}

%\author{First Author\affilnum{1} and Second Author\affilnum{2}}

%\affiliation{\affilnum{1}Department, Institution, Country\\
%\affilnum{2}Department, Institution, Country}

%\corrauth{Corresponding Author Name, Department, Institution, Address.}

\email{corresponding.author@email.edu}

\begin{abstract}
% Write a concise abstract here.
\end{abstract}

\keywords{keyword 1, keyword 2, keyword 3}

\maketitle

%=============================
%======== INTRODUCTION
%=============================
\section{Introduction}

% Introduce the problem, motivation, gap in knowledge, and study objective.

\begin{figure*}[ht]
    \centering
    %\includegraphics[width=0.999\textwidth]{figures/figure1}
    \caption{
    \textbf{Figure title.}
    Brief figure description.
    }
    \label{fig:figure1}
\end{figure*}

%=============================
%======== METHODS
%=============================
\section{Methods}

% Describe the overall study design, framework, model, experiment, or analysis pipeline.

\subsection{Study Design or Framework Overview}

% Describe the overall design or conceptual framework.

\subsection{Data Sources or Materials}

% Describe datasets, experimental materials, participants, simulations, or inputs.

\subsection{Model, Algorithm, or Analytical Method}

% Describe the core method, computational model, statistical model, or experimental procedure.

\subsection{Validation, Evaluation, or Statistical Analysis}

% Describe validation strategy, metrics, uncertainty analysis, and statistical tests.

\subsection{Implementation or Reproducibility}

% Describe software, code availability, computational environment, or reproducibility details.

%=============================
%======== RESULTS
%=============================
\section{Results}

% Present the main findings in a logical order.

\subsection{Primary Finding}

% Describe the first major result.

\subsection{Secondary Finding}

% Describe the second major result.

\subsection{Sensitivity, Robustness, or Subgroup Analysis}

% Describe additional analyses, robustness checks, or subgroup results.

\begin{table*}[ht]
\footnotesize\sf
\centering
\caption{Table title.\label{tab:table1}}
\begin{tabular}{lll}
\hline
\textbf{Column 1} & \textbf{Column 2} & \textbf{Column 3} \\
\hline
% Row 1 & Row 1 & Row 1 \\
% Row 2 & Row 2 & Row 2 \\
\hline
\end{tabular}
\end{table*}

%=============================
%======== DISCUSSION
%=============================
\section{Discussion}

% Interpret the findings, explain their significance, and relate them to prior work.

\subsection{Principal Findings}

% Summarize the most important findings without repeating all results.

\subsection{Comparison With Prior Work}

% Compare findings with existing literature.

\subsection{Implications}

% Discuss scientific, clinical, educational, engineering, or translational implications.

\subsection{Limitations}

% Discuss limitations, assumptions, uncertainties, and generalizability.

\subsection{Future Work}

% Describe next steps and future research directions.

%=============================
%======== CONCLUSION
%=============================
\section{Conclusion}

% Provide a brief concluding paragraph that reinforces the main contribution.

%=============================
%======== APPENDIX
%=============================
\section*{Appendix A. Supplementary Methods or Data Architecture}

% Add supplementary methodological details, schemas, derivations, or extended descriptions.

\subsection*{A.1 Data Sources or Inputs}

\begin{itemize}[leftmargin=*,itemsep=3pt]
    \item Placeholder item 1.
    \item Placeholder item 2.
    \item Placeholder item 3.
\end{itemize}

\subsection*{A.2 Processing or Ingestion Pipeline}

\begin{enumerate}[leftmargin=*,itemsep=3pt]
    \item \textbf{Step 1:} Description.
    \item \textbf{Step 2:} Description.
    \item \textbf{Step 3:} Description.
\end{enumerate}

\subsection*{A.3 Core Schemas, Variables, or Parameters}

\begin{itemize}[leftmargin=*,itemsep=3pt]
    \item \textbf{Category 1:} Description.
    \item \textbf{Category 2:} Description.
    \item \textbf{Category 3:} Description.
\end{itemize}

\section*{Appendix B. Supplementary Validation, Compliance, or Implementation Details}

% Add extended validation results, governance details, reproducibility notes, or additional tables.

\subsection*{B.1 Architecture or Implementation}

\begin{itemize}[leftmargin=*,itemsep=3pt]
    \item Placeholder item 1.
    \item Placeholder item 2.
    \item Placeholder item 3.
\end{itemize}

\subsection*{B.2 Evaluation or Compliance Considerations}

\begin{itemize}[leftmargin=*,itemsep=3pt]
    \item Placeholder item 1.
    \item Placeholder item 2.
    \item Placeholder item 3.
\end{itemize}

\subsection*{B.3 Risk, Safety, or Quality-Control Notes}

\begin{itemize}[leftmargin=*,itemsep=3pt]
    \item Placeholder item 1.
    \item Placeholder item 2.
    \item Placeholder item 3.
\end{itemize}

\bibliographystyle{plain}
\bibliography{source}

\end{document}
